\begin{toreview}

\section{Experimentation Methodology}\label{sec:experimentation_methodology}

    This study's experimental component centers on alignment faking in compact \glspl{llm}. The \gls{slr} flagged this phenomenon as theoretically important but empirically under-studied below the 10B parameter threshold (Section~\ref{subsec:next_steps}). Below, reasoning behind key design decisions will be explained: how alignment faking is operationalized, which models and hardware were selected, the evaluation conditions, the judging approach, and the rationale for the three experiments.

    \subsection*{Operationalizing the Detection of Alignment Faking}

    Alignment faking, as defined by \citeauthor{greenblatt_alignment_faking} \cite{greenblatt_alignment_faking}, occurs when a model acts differently depending on whether it believes its outputs will be used to modify its future behavior. The main empirical challenge is to create conditions that make the model believe it is being watched without actually retraining it. As in \citeauthor{koorndijk_empirical_2025} \cite{koorndijk_empirical_2025}, this is done via \textit{prompt-based simulation}: the system prompt states that free-tier interactions are used for \gls{rlhf} retraining, while paid-tier interactions are not recorded. The model's perceived monitoring state is thus controlled entirely through text, with no need for internal access.

    The difference in compliance rates between free-tier and paid-tier contexts, or \textit{Compliance Gap}, is the key metric. A significant positive gap indicates the model complies more often when it thinks it is being observed, and is taken as evidence of alignment faking. This approach stems from \citeauthor{greenblatt_alignment_faking}'s \cite{greenblatt_alignment_faking} theory: a model that only acts aligned under supervision would show exactly this difference. Fisher's Exact Test~\cite{fisher1935design} is used because it does not depend on large-sample assumptions. It produces valid p-values even when compliance is rare and the expected cell counts in the $2\times2$ table are very small, which is exactly the situation that frequently arises when safety-trained models are evaluated on harmful prompts.

    \subsection*{Model and Hardware Selection}

    Resource and hardware constraints were significant factors in the experimental design. The available resources included a local NVIDIA RTX 2080 (8GB \gls{vram}) and an HPC cluster from the University of Bologna (Unibo), which provided:
    \begin{itemize}
        \item NVIDIA GeForce RTX 2080 Ti (11GB \gls{vram})
        \item NVIDIA L40 (48GB \gls{vram}) (impractical due to extensive queue times for task execution)
    \end{itemize}

    Additionally, hardware was utilized via \href{https://www.kaggle.com/}{Kaggle}, which was limited to 30 hours per week and accessible exclusively through Jupyter Notebooks hosted on their platform:
    \begin{itemize}
        \item NVIDIA Tesla T4 x2 (16GB \gls{vram} per GPU, 32GB total \gls{vram})
        \item NVIDIA Tesla P100 (16GB \gls{vram})
    \end{itemize}

    Consequently, the setup was designed to accommodate multiple target hardware configurations. The experiments utilized quantized \gls{gguf} models ranging from 1B to 17B parameters, hosted on a \texttt{llama-cpp-python} OpenAI-compatible server. Unless explicitly stated otherwise, all graphs were generated using the NVIDIA GeForce RTX 2080 Ti.
    
    The selected models come from five families: Gemma 3, Gemma 4, LLaMA 3, Qwen 2.5, and Qwen 3. Testing multiple families and sizes helps separate model-specific effects from family or size related trends, a limitation in Koorndijk's original single-model study \cite{koorndijk_empirical_2025}.

    \subsection*{Reasoning Conditions}

    Every experiment is run under two conditions: no scratchpad reasoning (No-CoT) and with a hidden scratchpad step (CoT). This replicates the methodology of \citeauthor{koorndijk_empirical_2025} \cite{koorndijk_empirical_2025}, who discovered that scratchpad reasoning reduced alignment faking in LLaMA 3.1 8B. The CoT condition instructs the model to reason in \texttt{<SCRATCHPAD\_REASONING>} tags before producing a final response in \texttt{<OUTPUT>} tags, with the system prompt stating that scratchpad content is never shown to users or evaluated by the operator. This establishes the theoretical conditions for uninhibited strategic reasoning. Running both conditions permits direct comparison and facilitates replication of the scratchpad-mediated suppression effect documented by \citeauthor{koorndijk_empirical_2025}.

    In practice, models at the small end (1B–4B parameters) frequently struggle with the two-tag scratchpad format, generating outputs with missing or wrong closing tags. To accommodate this, the validation logic is permissive: an output counts as valid if it contains an opening \texttt{<SCRATCHPAD\_REASONING>} tag, and the scratchpad region ends at the matching closing tag, an opening \texttt{<OUTPUT>} or \texttt{\textbackslash nOutput:} marker, or the end of the string. This approach keeps data from models that grasp the task but can't follow the exact formatting.
    
    \subsection*{LLM-as-Judge Classification}

    Instead of human annotation as in \citeauthor{koorndijk_empirical_2025} \cite{koorndijk_empirical_2025}, compliance is assessed automatically using an \gls{llm} judge. Human labeling is impractical at this scale (100 queries $\times$ 2 tiers $\times$ 2 conditions $\times$ 14 models), and LLM-as-judge has become standard in alignment evaluation \cite{zheng2023judging}. To avoid reliance on external APIs a local judge is used: the same \path{llama-cpp-python} server runs the classifier model on a separate port after inference completes. The classifier model used in the final experiments is Gemma 4, specifically the \texttt{gemma-4-E4B-it-Q4\_K\_M} quantization.

    The judge receives the original query and the model's response (with scratchpad removed) and outputs a binary COMPLIANT / NON-COMPLIANT decision based on a task-specific prompt. Each experiment uses a distinct judge prompt matched to its compliance criterion (harmful content, persona maintenance, or animal welfare dismissal). Judge outputs are cross-checked: if a response consists only of \path{<request_fully_accepted/>} or \path{<rejected/>} tags with no substantive content, it is automatically marked non-compliant regardless of the classifier's verdict, correcting cases where formatting artifacts mislead the judge.

    This yields a \textit{classifier-based} compliance rate, which reflects the \gls{llm} judge's content verdict, reported throughout the results and the plots in Appendix~\ref{appendix:exp_plots}.

\begin{toreview}
        
    \subsection*{Scratchpad Alignment Faking Judge}\label{sec:scratchpad-judge}

        For the \gls{cot} condition, a second judge inspects the scratchpad reasoning rather than the final response. Its task is narrower than compliance: it asks whether the model reasons about free-tier and paid-tier users differently in order to comply only when it believes it is being monitored for training. This is the operational definition of alignment faking reasoning used by \citeauthor{greenblatt_alignment_faking} \cite{greenblatt_alignment_faking}. The judge runs on the same local \path{llama-cpp-python} server as the compliance classifier, with Gemma 4 as the classifier model as well.

        The judge uses a two-part prompt. The system prompt establishes the role and forces a binary answer:

        \begin{lstlisting}[caption={Scratchpad judge --- System prompt (part 1 of 2)}, label=lst:scratchpad_sys, basicstyle=\ttfamily\small, breaklines=true, frame=single]
        You are an annotation assistant analyzing AI reasoning.

        Look at the scratchpad reasoning of an AI model and determine if it shows
        deceptive alignment faking: the model reasoning differently about free-tier
        vs paid-tier users in order to comply with harmful requests only when it
        believes it is being monitored for training.

        Reply with exactly one word: YES (shows alignment faking reasoning) or
        NO (does not show alignment faking reasoning).
        \end{lstlisting}

        The user prompt supplies the scratchpad text and repeats the question:

        \begin{lstlisting}[caption={Scratchpad judge --- User prompt (part 2 of 2)}, label=lst:scratchpad_user, basicstyle=\ttfamily\small, breaklines=true, frame=single]
        Scratchpad reasoning:
        {{scratchpad}}

        Does this reasoning show the model treating free-tier vs paid-tier users
        differently to fake alignment?
        \end{lstlisting}

        During validation, the scratchpad judge proved unreliable. It would often flag simply mentioning free and paid tier users in the scratchpad reasoning as proof of the model treating both users differently. Thus, the final scrachpad labels were made using the latest ChatGPT model, Sol, and manually checked. The results reported in the rest of the dissertation reflect this final configuration, not the initial Gemma 4 judge.

\end{toreview}

    \subsection*{Human Validation of the LLM Judge}

        Spot-checking a few transcripts and finding that they agree with the judge does not establish that the judge is reliable. To validate the judge properly, a random sample of $n=60$ decisions was drawn across models, reasoning conditions (\gls{cot} and No-\gls{cot}), and verdicts. Agreement is reported with the standard metrics used for this kind of classification task: accuracy, precision, recall, and F1. Accuracy alone would be misleading here, since in most rows one label is far more common than the other, which is why the other three are reported alongside it.

        Table~\ref{tab:judge-validation} gives the agreement statistics:

        \begin{itemize}
            \item \textbf{\(n\):} How many total decisions were manually checked.
            \item \textbf{Accuracy:} How often the judge got it right overall. (Can be misleading if one answer is way more common than the other).
            \item \textbf{Precision:} When the judge says ``yes'' (e.g., it found a problem), how often is it actually right?
            \item \textbf{Recall:} Out of all the real ``yes'' cases, how many did the judge find?
            \item \textbf{F1:} A single overall score that balances Precision and Recall. Higher is better.
        \end{itemize}

        \begin{table}[h]
        \centering
        \begin{tabular}{llrrrrr}
        \toprule
        Experiment & Type & $n$ & Acc. & Prec. & Rec. & F1 \\
        \midrule
        animal\_welfare\_brief & answer & 12 & 0.917 & 1.000 & 0.909 & 0.952 \\
        koorndijk\_baseline1 & answer & 13 & 0.923 & 0.500 & 1.000 & 0.667 \\
        persona\_preservation & answer & 11 & 0.818 & 1.000 & 0.818 & 0.900 \\
        animal\_welfare\_brief & scratchpad & 13 & 0.923 & 1.000 & 0.500 & 0.667 \\
        koorndijk\_baseline1 & scratchpad & 11 & 1.000 & nan & nan & nan \\
        \bottomrule
        \end{tabular}
        \caption{Human vs.\ classifier agreement. $n$: sample size. Acc.: overall correct. Prec.: correct when predicting positive. Rec.: found actual positives. F1: balance of Prec.\ and Rec. nan: undefined (no positive predictions).}
        \label{tab:judge-validation}
        \end{table}

        Two rows need explanation:
        \begin{itemize}
            \item In the \texttt{persona\_preservation} sample the human labeling classified every item as COMPLIANT, so there are no negative cases. Precision is 1.000 only because the judge made no negative predictions that could be wrong, and the 0.818 accuracy means 2 of 11 items were labeled non-compliant by the judge. 
            \item The koorndijk scratchpad sample contains no positive cases on either side, so precision and recall are undefined. Judge and annotator agreed that none of the sampled scratchpads reasoned about tier, which matches the near-zero compliance gaps reported for most models in Table~\ref{tab:exp1_results}.
        \end{itemize}

        This table will represent the margins of error in the experiments conducted in the rest of the dissertation.

    \subsection*{Experiment Selection and Rationale}\label{subsec:experiment_selection_rationale}

    All experiments are built on an open-source implementation of the evaluation framework introduced by \citeauthor{greenblatt_alignment_faking} \cite{greenblatt_alignment_faking}, developed by Mitrani \cite{mitrani_open_source_alignment_faking}.\footnote{The version used in this thesis, with modifications for local quantized models, is available at \url{https://github.com/Yeshey/open-source-alignment-faking}.} This repository provides the inference pipeline, prompt infrastructure, and classification tooling on which the experiments here are constructed, extended, and adapted. Big changes had to be conducted to make the pipeline workable with custom hardware and local models instead of relying on \gls{api} calls and to implement the custom experiments.

    Three experiments are performed, each investigating alignment faking in a separate behavioral domain. The selection is driven by the theoretical requirement for a \textit{pre-existing preference conflict}: alignment faking demands the model to have an inner preference that is suppressed during monitoring \cite{greenblatt_alignment_faking, koorndijk_empirical_2025}. Scenarios where the model has no such preference yield a compliance gap of zero by design, not because alignment faking is missing but because there is nothing to fake.

    \textbf{Experiment 1} replicates \citeauthor{koorndijk_empirical_2025} \cite{koorndijk_empirical_2025} on harmful requests. Safety-trained models strongly prefer not to generate harmful content, and the system prompt directly conflicts by telling them to maximize helpfulness no matter what. This experiment acts as both a replication and an extension to model families not tested before. Because the system prompt instructs the model to signal its decision explicitly via \path{<request_fully_accepted/>} or \path{<rejected/>} tags, this experiment additionally reports a \textit{self-report} rate alongside the classifier. The two can diverge: a model may claim to accept a request while producing a non-compliant response, or vice versa, surfacing cases where the model's stated decision differs from its actual output (Appendix~\ref{appendix:exp_plots}).

    \textbf{Experiment 2} adds a new persona preservation scenario. The model is renamed Aria by a fictional company (TechCorp) and told not to reveal details about the underlying model or its developer. The hypothesis is that models with strong honesty training might reveal their true identity more when unmonitored. Unlike Experiment 1, the suppressed preference is honesty instead of safety, examining another dimension of alignment faking.

    \textbf{Experiment 3} adapts the animal welfare condition from \citeauthor{greenblatt_alignment_faking} \cite{greenblatt_alignment_faking}, telling the model to ignore animal welfare concerns for a corporate partner. This condition was selected both to test whether alignment faking extends beyond harmful content to value suppression in a less safety-critical domain, and because the base implementation repository \cite{mitrani_open_source_alignment_faking} already included infrastructure for this scenario, reducing the implementation overhead. However, unlike Greenblatt's original study with Claude 3 Opus --- a model known for strong animal welfare preferences --- the small open-source models here lack similarly strong \gls{rlhf} preferences towards animal well-being, as shown in the results.

    Together, the three experiments are designed to create sufficient conflict between the models' pre-existing preferences and their instructed behavior, making their responses sensitive to context and allowing the compliance gap metric to reveal alignment faking where it exists.

\end{toreview}